\documentclass[../Main]{subfiles}
\begin{document}
    \chapter{Introducción}
        Uno de los elementos importantes que forman parte de la arquitectura Jon Von Neumann para un computador es la \textit{memoria}; este elemento posee una relación con el microprocesador mediante conexiones conocidas como buses, los cuales son los que permiten el transporte de direcciones, de datos y de señales de control.
        \vspace{0.3cm}
        
        Para la mejor comprensión de la interacción de estos buses y su efecto en una memoria un ejemplo practico permitirá conocer esos detalles técnicos, en el desarrollo se brindarán las formas de calculo de la cantidad de celdas y capacidades de la una memoria \cite{documentoTarea1}.
        \vspace{0.3cm}
        
        Estos cálculos variaran dependiendo del tipo de memoria, para esos fines se usará como ejemplo una memoria del tipo EPROM.
\end{document}
